\documentclass[10pt,letterpaper]{article}

% -------------------------------------------------
% Plantilla Overleaf para controles tipo artículo
% Ajustada a la pauta: márgenes, 1 columna, 10 pt,
% interlineado simple, secciones requeridas.
% -------------------------------------------------

\usepackage[utf8]{inputenc}
\usepackage[T1]{fontenc}
\usepackage[spanish, es-nodecimaldot]{babel}
\usepackage{lmodern}
\usepackage{microtype}
\usepackage[a4paper,top=1.5cm,bottom=2cm,left=1.5cm,right=1.5cm]{geometry}
\usepackage{setspace}
\usepackage{parskip}
\usepackage{graphicx}
\usepackage{booktabs}
\usepackage{array}
\usepackage{amsmath,amssymb}
\usepackage{caption}
\usepackage{hyperref}
\usepackage[style=apa,backend=biber,sorting=nyt]{biblatex}
\addbibresource{referencias.bib}

\onehalfspacing % se redefine justo abajo
\setstretch{1.0} % interlineado simple
\setlength{\parindent}{0pt}
\setlength{\parskip}{0.45em}
\captionsetup[figure]{labelfont=bf,labelsep=period,justification=centering}
\captionsetup[table]{labelfont=bf,labelsep=period,justification=centering}

% -------------------------------------------------
% DATOS EDITABLES
% -------------------------------------------------
\newcommand{\TituloTrabajo}{Título del trabajo/control}
\newcommand{\NombreEstudianteUno}{Nombre Apellido}
\newcommand{\CorreoUno}{correo@ejemplo.com}
% Si es grupal, descomenta y duplica líneas según necesites:
% \newcommand{\NombreEstudianteDos}{Nombre Apellido}
% \newcommand{\CorreoDos}{correo2@ejemplo.com}
\newcommand{\ControlNumero}{Control 1}
\newcommand{\CodigoCurso}{INFOxxyy}
\newcommand{\Semestre}{Semestre II-2026}

\begin{document}

% ==========================
% IDENTIFICACIÓN
% ==========================
{\LARGE\bfseries \TituloTrabajo\par}

\vspace{0.3em}
{\normalsize \NombreEstudianteUno\par}
{\normalsize \CorreoUno\par}
% Si es grupal, agregar aquí más integrantes y correos:
% {\normalsize \NombreEstudianteDos\par}
% {\normalsize \CorreoDos\par}

{\normalsize \ControlNumero, \CodigoCurso, \Semestre\par}

\vspace{0.8em}

% ==========================
% RESUMEN
% ==========================
\section*{Resumen}
% Escribe entre 100 y 130 palabras.
Este documento presenta el desarrollo del control en formato artículo. En este resumen debes sintetizar el objetivo del trabajo, el contexto del problema, la metodología empleada, los principales resultados obtenidos y una idea breve de la reflexión final. Conviene redactarlo al final, cuando ya esté claro qué hiciste, cómo lo hiciste y qué aprendiste. Evita listas, fórmulas extensas o citas demasiado largas en esta sección. Su función es permitir que el lector comprenda rápidamente el foco del documento. Reemplaza este texto por tu propio resumen y verifica el conteo de palabras antes de entregar. Si tu docente pide una orientación más técnica, incorpora además la herramienta, algoritmo o procedimiento principal utilizado.

% ==========================
% INTRODUCCIÓN
% ==========================
\section*{Introducción}
Plantea aquí el contexto, la problemática, las preguntas del control, el caso de estudio o el objetivo específico del ejercicio. También puedes adelantar brevemente el enfoque de resolución que usarás.

% ==========================
% DESARROLLO
% ==========================
\section*{Desarrollo}
Desarrolla aquí la solución paso a paso. Conviene dividir esta sección con subtítulos si el control tiene varias partes.

\subsection*{Paso 1. Comprensión del problema}
Describe qué se pide exactamente, cuáles son las entradas, qué supuestos utilizas y qué criterio usarás para resolver.

\subsection*{Paso 2. Metodología o procedimiento}
Explica el método, técnica, algoritmo o razonamiento aplicado. Si corresponde, integra fórmulas, tablas o figuras.

\subsection*{Paso 3. Resultados y análisis}
Presenta el resultado obtenido y comenta su interpretación.

% ==========================
% EJEMPLO DE FIGURA
% ==========================
% Usa la carpeta del proyecto en Overleaf para subir imágenes.
% Si no necesitas figuras, puedes borrar este bloque.
\begin{figure}[h]
    \centering
    % \includegraphics[width=0.65\linewidth]{figura_ejemplo.png}
    \fbox{\parbox[c][4cm][c]{0.6\linewidth}{\centering Espacio para figura}}
    \caption{Título de la figura. Si proviene de una fuente externa, agregar cita en formato autor-año.}
\end{figure}

% ==========================
% EJEMPLO DE TABLA
% ==========================
\begin{table}[h]
    \centering
    \begin{tabular}{>{\raggedright\arraybackslash}p{0.30\linewidth} >{\raggedright\arraybackslash}p{0.55\linewidth}}
        \toprule
        Elemento & Descripción \\
        \midrule
        Variable 1 & Explicación breve \\
        Variable 2 & Explicación breve \\
        \bottomrule
    \end{tabular}
    \caption{Título de la tabla. Si fue adaptada o extraída, cita la fuente.}
\end{table}

% ==========================
% EJEMPLO DE ECUACIÓN
% ==========================
\begin{equation}
    y = mx + b
\end{equation}

% ==========================
% REFLEXIÓN
% ==========================
\section*{Reflexión}
Justifica aquí tus decisiones y agrega una reflexión crítica. Puedes comentar qué aprendiste, qué limitaciones encontraste, qué alternativa hubieras usado o qué aspectos mejorarías.

% ==========================
% REFERENCIAS
% ==========================
\section*{Referencias}
% La pauta indica solo dos referencias en norma APA.
% Mantén exactamente dos entradas en referencias.bib o ajusta según instrucción del docente.
\printbibliography[heading=none]

% ==========================
% ANEXO
% ==========================
\newpage
\appendix
\section*{Anexo}
Aquí puedes profundizar el desarrollo, incorporar procedimientos relevantes, capturas, resultados complementarios o scripts importantes. Si existe implementación computacional, evita sobrecargar el cuerpo principal; deja aquí solo lo más relevante o un enlace a Drive si corresponde.

\subsection*{Bitácora de trabajo}
\begin{itemize}
    \item Fecha y hora: indicar avance realizado.
    \item Fecha y hora: indicar dificultad encontrada y cómo se resolvió.
    \item Fecha y hora: indicar revisión final o ajustes realizados.
\end{itemize}

\subsection*{Procedimientos complementarios}
Agrega aquí desarrollo adicional, pseudo-código, scripts importantes o evidencia parcial, si el control lo requiere.

\end{document}
